
\documentclass{article}
\usepackage{colortbl}
\usepackage{makecell}
\usepackage{multirow}
\usepackage{supertabular}

\begin{document}

\newcounter{utterance}

\centering \large Interaction Transcript for game `cladder', experiment `full\_v1.5\_default', episode 188 with llama{-}3.1{-}8b{-}instruct.
\vspace{24pt}

{ \footnotesize  \setcounter{utterance}{1}
\setlength{\tabcolsep}{0pt}
\begin{supertabular}{c@{$\;$}|p{.15\linewidth}@{}p{.15\linewidth}p{.15\linewidth}p{.15\linewidth}p{.15\linewidth}p{.15\linewidth}}
    \# & \multicolumn{2}{c}{Player} && \multicolumn{2}{c}{Game Master} \\
    \hline

    \theutterance \stepcounter{utterance}  
    & & & \multicolumn{4}{p{0.6\linewidth}}{
        \cellcolor[rgb]{0.9,0.9,0.9}{
            \makecell[{{p{\linewidth}}}]{
                \texttt{\tiny{[P1$\langle$GM]}}
                \texttt{You are an expert in causal inference. The following question is not a typical commonsense query, but rather a meticulously designed question created by a professor specializing in causal inference, intended to assess the students' mastery of the course content.} \\
\texttt{Imagine a self{-}contained, hypothetical world with only the following conditions, and without any unmentioned factors or causal relationships: The captain has a direct effect on the private and the corporal. The corporal has a direct effect on prisoner. The private has a direct effect on prisoner. Method 1: We look directly at how the captain correlates with prisoner in general. Method 2: We look at this correlation case by case according to the private. To understand how the captain affects prisoner, is it more correct to use the Method 1 than Method 2?} \\
\texttt{Start your answer with 'Yes' or 'No', followed by additional reasoning or evidence to support your explanation.} \\
            }
        }
    }
     \\ \\

    \theutterance \stepcounter{utterance}  
    & \multicolumn{4}{p{0.6\linewidth}}{
        \cellcolor[rgb]{0.561,0.737,0.561}{
            \makecell[{{p{\linewidth}}}]{
                \texttt{\tiny{[P1$\rangle$GM]}}
                \texttt{No, it is not more correct to use Method 1 than Method 2 to understand how the captain affects prisoner.} \\
\\ 
\texttt{In this hypothetical world, the captain has a direct effect on both the private and the corporal, who in turn have direct effects on the prisoner. This creates a situation where the captain's effect on the prisoner is mediated by both the private and the corporal.} \\
\\ 
\texttt{Method 1, which looks directly at the correlation between the captain and the prisoner in general, would likely overestimate the direct effect of the captain on the prisoner, as it does not account for the indirect effects mediated by the private and the corporal. This is because the correlation between the captain and the prisoner would capture both the direct and indirect effects, leading to an inflated estimate of the captain's causal effect.} \\
\\ 
\texttt{On the other hand, Method 2, which looks at the correlation case by case according to the private, would allow us to control for the private's effect on the prisoner and estimate the captain's effect on the prisoner while holding the private's effect constant. This would provide a more accurate estimate of the captain's direct effect on the prisoner, as it would isolate the effect of the captain from the indirect effects mediated by the private.} \\
\\ 
\texttt{Therefore, Method 2 is a more correct approach to understanding how the captain affects the prisoner, as it allows for a more precise estimation of the captain's causal effect while controlling for the mediating effects of the private and the corporal.} \\
            }
        }
    }
    & & \\ \\

    \theutterance \stepcounter{utterance}  
    & & & \multicolumn{2}{p{0.3\linewidth}}{
        \cellcolor[rgb]{0.9,0.9,0.9}{
            \makecell[{{p{\linewidth}}}]{
                \texttt{\tiny{[GM$|$GM]}}
                \texttt{No} \\
            }
        }
    }
    & & \\ \\

    \theutterance \stepcounter{utterance}  
    & & & \multicolumn{2}{p{0.3\linewidth}}{
        \cellcolor[rgb]{0.9,0.9,0.9}{
            \makecell[{{p{\linewidth}}}]{
                \texttt{\tiny{[GM$|$GM]}}
                \texttt{yes} \\
            }
        }
    }
    & & \\ \\

    \theutterance \stepcounter{utterance}  
    & & & \multicolumn{2}{p{0.3\linewidth}}{
        \cellcolor[rgb]{0.9,0.9,0.9}{
            \makecell[{{p{\linewidth}}}]{
                \texttt{\tiny{[GM$|$GM]}}
                \texttt{game\_result = LOSE} \\
            }
        }
    }
    & & \\ \\

\end{supertabular}
}

\end{document}
